\section{Requirements}\label{sec:requirements}
This section defines the functional and non-functional requirements.

\subsection{Functional Requirements}\label{subsec:functional-requirements}
These requirements specify how the system implements core systems, AI tooling, and expected behavior from users.

\subsubsection{Framework Requirements}
\begin{itemize}
    \item FR1: The framework must provide a deterministic routing system.
    \item FR2: Routing system must support common HTTP methods.
    \item FR3: Routing system must be modular, such that part of the routing tree can be moved to different branch
        and all the bindings must work as expected.
    \item FR4: The framework has the ability to automatically infer CRUD operations for database tables.
    \item FR5: Models must map database columns to properties without extra boilerplate.
    \item FR6: The framework must provide form-handling subsystem.
    \item FR7: The system has the ability to automatically infer forms based on model structure.
    \item FR8: The framework must provide a component-based editor capable of creating and maintaining webpages.
    \item FR9: Editor components must be rendered on client from provided structure.
\end{itemize}

\subsubsection{AI Generation Requirements}
\begin{itemize}
    \item AIR1: The system must include an AI-powered translation unit capable of translating static text and text templates for supported languages.
    \item AIR2: The framework must support generating new editor components.
    \item AIR3: Users must be able to request AI-generated static webpages from natural language prompt.
    \item AIR4: The system must provide AI-documentation from its own source files.
    \item AIR5: The system must include website structure generation of pages and their templates.
\end{itemize}

\subsection{Non-Functional Requirements}\label{subsec:non-functional-requirements}
These requirements specify how the whole system behaves rather than addressing specific functionalities.

\begin{itemize}
    \item NFR1: The system must prioritize ease of development, readability, and extensibility.
    \item NFR2: New features, components, and templates must be easy to add and maintain.
    \item NFR3: The framework must perform efficiently enough for small-to-medium projects.
    \item NFR4: All core features must be implemented with minimal reliance on third-party libraries.
    \item NFR5: The code must follow the conventions defined by the framework’s philosophy.
    \item NFR6: The framework's systems must follow a unified design approach.
\end{itemize}
